\documentclass[11pt]{article}

\usepackage[margin=0.8in]{geometry}
\usepackage{amsmath,amssymb}
\usepackage{booktabs}
\usepackage{array}
\usepackage{enumitem}
\usepackage[hidelinks]{hyperref}

\setlist[itemize]{leftmargin=1.4em}
\setlist[enumerate]{leftmargin=1.6em}

\title{SYSC 5108 Exam Study\\W25 Sample Final Step-by-Step Solutions}
\author{Jesse Levine}
\date{}

\begin{document}
\maketitle

\section*{Source Note}

This study note is based on the two files in the course directory:
\begin{itemize}
    \item \texttt{W25\_SYSC5108\_Final.pdf}
    \item \texttt{W25\_SYSC5108\_Final\_Solution.pdf}
\end{itemize}

The PDF filename says ``Final'', while the cover page itself says ``Midterm I''.
I follow the file naming used in the directory and turn the provided sample exam
and sample solution into a cleaner step-by-step study document.

\section*{Question 1: Linear Regression}

Training points:
\[
    (0,-1),\ (1,2),\ (6,4),\ (9,2)
\]
Test point:
\[
    (3,1).
\]

We fit the line
\[
    \hat y = \beta_0 + \beta_1 x.
\]

\subsection*{Step 1: Build the design matrix}

\[
    X =
    \begin{bmatrix}
        1 & 0 \\
        1 & 1 \\
        1 & 6 \\
        1 & 9
    \end{bmatrix},
    \qquad
    y =
    \begin{bmatrix}
        -1 \\ 2 \\ 4 \\ 2
    \end{bmatrix}.
\]

\subsection*{Step 2: Use the normal equation}

\[
    \beta = (X^\top X)^{-1}X^\top y.
\]

Compute
\[
    X^\top X =
    \begin{bmatrix}
        4 & 16 \\
        16 & 118
    \end{bmatrix},
    \qquad
    X^\top y =
    \begin{bmatrix}
        7 \\ 44
    \end{bmatrix}.
\]

The determinant is
\[
    4(118)-16(16)=216.
\]

So
\[
    (X^\top X)^{-1}
    =
    \frac{1}{216}
    \begin{bmatrix}
        118 & -16 \\
        -16 & 4
    \end{bmatrix}.
\]

Then
\[
    \beta
    =
    \frac{1}{216}
    \begin{bmatrix}
        118 & -16 \\
        -16 & 4
    \end{bmatrix}
    \begin{bmatrix}
        7 \\ 44
    \end{bmatrix}
    =
    \begin{bmatrix}
        61/108 \\
        8/27
    \end{bmatrix}
    \approx
    \begin{bmatrix}
        0.5648 \\
        0.2963
    \end{bmatrix}.
\]

Therefore the fitted model is
\[
    \boxed{\hat y = 0.5648 + 0.2963x.}
\]

\subsection*{Step 3: Predict the training outputs}

\[
    \hat y(0)=0.5648,\quad
    \hat y(1)=0.8611,\quad
    \hat y(6)=2.3426,\quad
    \hat y(9)=3.2315.
\]

\subsection*{Step 4: Compute the training MSE}

\[
    \text{MSE}_{\text{train}}
    =
    \frac{1}{4}\sum_{i=1}^{4}(y_i-\hat y_i)^2.
\]

The squared errors are approximately
\[
    2.4481,\quad 1.2971,\quad 2.7467,\quad 1.5165.
\]

So
\[
    \text{MSE}_{\text{train}}
    \approx
    \frac{2.4481+1.2971+2.7467+1.5165}{4}
    \approx 2.0021.
\]

Rounded for exam work:
\[
    \boxed{\text{MSE}_{\text{train}} \approx 2.0.}
\]

\subsection*{Step 5: Compute the test error}

For the test input \(x=3\),
\[
    \hat y(3)=0.5648+0.2963(3)\approx 1.4537.
\]

The squared test error is
\[
    (1-1.4537)^2 \approx 0.2058.
\]

Since there is only one test point, its MSE is the same number:
\[
    \boxed{\text{MSE}_{\text{test}} \approx 0.206.}
\]

\section*{Question 2: Multiclass Logistic Regression Derivative}

The exam gives
\[
    p(y=k\mid x,\beta)=\text{softmax}(z)_k
    =
    \frac{e^{z_k}}{\sum_j e^{z_j}},
\]
with
\[
    z_k = \beta_{k0} + \beta_{k1}x,
\]
and
\[
    J(\beta)
    =
    -\sum_{i=1}^{M}\sum_{k=1}^{K} I_k(y^{(i)})\log \text{softmax}(z^{(i)})_k.
\]

The target identity in the sample solution is
\[
    \nabla_{\beta_k}J(\beta_k)
    =
    -\sum_{i=1}^{M}\sum_{k=1}^{K}
    I_k(y^{(i)})
    \left(1-\text{softmax}(z^{(i)})_k\right)x^{(i)}.
\]

\subsection*{Step 1: Move the derivative inside the sum}

\[
    \nabla_{\beta_k}J(\beta_k)
    =
    -\sum_{i=1}^{M}\sum_{k=1}^{K}
    I_k(y^{(i)})
    \nabla_{\beta_k}\log \text{softmax}(z^{(i)})_k.
\]

\subsection*{Step 2: Differentiate the log}

Use
\[
    \nabla_{\beta_k}\log \text{softmax}(z)_k
    =
    \frac{\nabla_{\beta_k}\text{softmax}(z)_k}{\text{softmax}(z)_k}.
\]

\subsection*{Step 3: Differentiate the softmax term}

\[
    \text{softmax}(z)_k
    =
    \frac{e^{z_k}}{\sum_j e^{z_j}}.
\]

Using the quotient rule and the fact that
\[
    \nabla_{\beta_k} z_k = x,
\]
we get
\[
    \nabla_{\beta_k}\text{softmax}(z)_k
    =
    x\,
    \frac{e^{z_k}\left(\sum_j e^{z_j}\right)-e^{z_k}e^{z_k}}
         {\left(\sum_j e^{z_j}\right)^2}.
\]

Factor out the softmax:
\[
    \nabla_{\beta_k}\text{softmax}(z)_k
    =
    x\,\text{softmax}(z)_k\left(1-\text{softmax}(z)_k\right).
\]

\subsection*{Step 4: Divide by \(\text{softmax}(z)_k\)}

\[
    \nabla_{\beta_k}\log \text{softmax}(z)_k
    =
    x\left(1-\text{softmax}(z)_k\right).
\]

\subsection*{Step 5: Substitute back into the loss gradient}

\[
    \nabla_{\beta_k}J(\beta_k)
    =
    -\sum_{i=1}^{M}\sum_{k=1}^{K}
    I_k(y^{(i)})
    x^{(i)}
    \left(1-\text{softmax}(z^{(i)})_k\right).
\]

So the exam's target expression is
\[
    \boxed{
    \nabla_{\beta_k}J(\beta_k)
    =
    -\sum_{i=1}^{M}\sum_{k=1}^{K}
    I_k(y^{(i)})
    \left(1-\text{softmax}(z^{(i)})_k\right)x^{(i)}
    }.
\]

\subsection*{Study note}

The sample solution differentiates the \(k\)-th softmax term with respect to the
matching class parameter \(\beta_k\). In standard multiclass logistic regression,
the full familiar gradient is usually written as
\[
    \sum_{i=1}^{M}\left(p_k^{(i)}-I_k(y^{(i)})\right)x^{(i)}.
\]
For exam study, it is useful to recognize both forms.

\section*{Question 3: 1D CNN Backpropagation}

From the figure:
\begin{itemize}
    \item input \(=[1,1,1,0,0,0]\),
    \item convolution bias \(w_0=0.75\),
    \item shared kernel weights \(w_1=+1,\ w_2=+1,\ w_3=-1\),
    \item output-layer weights \(w_{70}=0.1,\ w_{75}=0.2,\ w_{76}=0.5\),
    \item ReLU activations everywhere,
    \item target output \(t=2\).
\end{itemize}

\subsection*{Step 1: Forward pass through the convolution layer}

Window 1 uses \((1,1,1)\):
\[
    z_1 = 0.75 + 1(1) + 1(1) - 1(1) = 1.75,
    \qquad
    a_1=\text{ReLU}(z_1)=1.75.
\]

Window 2 uses \((1,1,0)\):
\[
    z_2 = 0.75 + 1(1) + 1(1) - 1(0) = 2.75,
    \qquad
    a_2=2.75.
\]

Window 3 uses \((1,0,0)\):
\[
    z_3 = 0.75 + 1(1) + 1(0) - 1(0) = 1.75,
    \qquad
    a_3=1.75.
\]

Window 4 uses \((0,0,0)\):
\[
    z_4 = 0.75,
    \qquad
    a_4=0.75.
\]

\subsection*{Step 2: Max pooling}

\[
    a_5 = \max(a_1,a_2)=\max(1.75,2.75)=2.75,
\]
\[
    a_6 = \max(a_3,a_4)=\max(1.75,0.75)=1.75.
\]

So the winning conv units are neuron 2 for pooling unit 5 and neuron 3 for
pooling unit 6.

\subsection*{Step 3: Output neuron}

\[
    z_7 = w_{70} + w_{75}a_5 + w_{76}a_6
    = 0.1 + 0.2(2.75) + 0.5(1.75).
\]

\[
    z_7 = 0.1 + 0.55 + 0.875 = 1.525.
\]

Since ReLU is active,
\[
    a_7 = 1.525.
\]

\subsection*{Part (a): Find the \(\delta\) values}

Using squared error \(L=\tfrac12(t-a_7)^2\),
\[
    \delta_7 = \frac{\partial L}{\partial z_7}
    = a_7 - t = 1.525 - 2 = -0.475.
\]

Backpropagate to the pooling units:
\[
    \delta_5 = \delta_7 w_{75} = -0.475(0.2) = -0.095,
\]
\[
    \delta_6 = \delta_7 w_{76} = -0.475(0.5) = -0.2375.
\]

Because max pooling routes the gradient only to the winning input:
\[
    \delta_1 = 0,\qquad
    \delta_2 = -0.095,\qquad
    \delta_3 = -0.2375,\qquad
    \delta_4 = 0.
\]

Therefore
\[
    \boxed{
    \delta_1=0,\ 
    \delta_2=-0.095,\ 
    \delta_3=-0.2375,\ 
    \delta_4=0,\ 
    \delta_5=-0.095,\ 
    \delta_6=-0.2375,\ 
    \delta_7=-0.475
    }.
\]

\subsection*{Part (b): Weight updates for \(w_{75}\) and \(w_1\)}

For the fully connected weight:
\[
    \frac{\partial L}{\partial w_{75}}
    = \delta_7 a_5
    = -0.475(2.75)
    = -1.30625.
\]

So
\[
    \boxed{\frac{\partial L}{\partial w_{75}}=-1.30625\approx -1.31.}
\]

For the shared convolution weight \(w_1\), add the contributions from each
window:
\[
    \frac{\partial L}{\partial w_1}
    =
    \delta_1 x_1 + \delta_2 x_2 + \delta_3 x_3 + \delta_4 x_4.
\]

The corresponding first inputs in the four windows are \(1,1,1,0\), so
\[
    \frac{\partial L}{\partial w_1}
    =
    0(1) + (-0.095)(1) + (-0.2375)(1) + 0(0)
    = -0.3325.
\]

Thus
\[
    \boxed{\frac{\partial L}{\partial w_1}=-0.3325.}
\]

\section*{Question 4: Discounted Return}

Given \(\gamma=0.9\) and reward sequence
\[
    R_1 = 2,\qquad R_2=R_3=\cdots = 7,
\]
find \(G_0\) and \(G_1\).

\subsection*{Step 1: Compute \(G_1\)}

\[
    G_1 = R_2 + \gamma R_3 + \gamma^2 R_4 + \cdots
    = 7 + 0.9(7) + 0.9^2(7) + \cdots.
\]

This is a geometric series:
\[
    G_1 = 7\sum_{n=0}^{\infty}0.9^n = \frac{7}{1-0.9}=70.
\]

\subsection*{Step 2: Compute \(G_0\)}

\[
    G_0 = R_1 + \gamma G_1 = 2 + 0.9(70)=2+63=65.
\]

So
\[
    \boxed{G_1=70,\qquad G_0=65.}
\]

\section*{Question 5: Reinforcement Learning}

The environment is
\[
\begin{array}{ccccc}
\toprule
\text{State} & \text{Action} & \text{Next state} & \text{Probability} & \text{Reward} \\
\midrule
0 & 0 & 0 & 1 & -1 \\
0 & 1 & 1 & 1 & 1 \\
1 & 0 & 1 & 1 & 1 \\
1 & 1 & 0 & 1 & -1 \\
\bottomrule
\end{array}
\]
\[
    \gamma=0.1,
    \qquad
    \pi(1\mid 0)=\pi(0\mid 1)=0.6,
    \qquad
    \alpha=0.6.
\]

\subsection*{Part (a): Generate the trajectories}

The random numbers are
\[
    0.94,\ 0.27,\ 0.30,\ 0.18,\ 0.75,\ 0.35.
\]

Following the same action-selection convention as the sample solution:
\begin{itemize}
    \item in state \(0\), choose action \(1\) with probability \(0.6\),
    \item in state \(1\), choose action \(0\) with probability \(0.6\).
\end{itemize}

\subsubsection*{Trajectory from state 0}

\begin{itemize}
    \item \(0.94 > 0.4\), so choose action \(1\): \(0 \to 1\), reward \(1\).
    \item \(0.27 < 0.6\), so in state \(1\) choose action \(0\): \(1 \to 1\), reward \(1\).
    \item \(0.30 < 0.6\), so again choose action \(0\): \(1 \to 1\), reward \(1\).
\end{itemize}

Thus
\[
    \tau_0: 0 \to 1 \to 1 \to 1,
    \qquad
    r = (1,1,1).
\]

\subsubsection*{Trajectory from state 1}

\begin{itemize}
    \item \(0.18 < 0.6\), so choose action \(0\): \(1 \to 1\), reward \(1\).
    \item \(0.75 > 0.6\), so choose action \(1\): \(1 \to 0\), reward \(-1\).
    \item \(0.35 < 0.4\), so in state \(0\) choose action \(0\): \(0 \to 0\), reward \(-1\).
\end{itemize}

Thus
\[
    \tau_1: 1 \to 1 \to 0 \to 0,
    \qquad
    r = (1,-1,-1).
\]

\subsection*{Monte Carlo values following the provided sample solution}

The sample solution updates to the following values:
\[
    v(0)\approx -0.156,
    \qquad
    v(1)\approx 0.974.
\]

For the action values:
\[
    q(0,0)=-1,
    \qquad
    q(0,1)=1+0.1(1)+0.1^2(1)=1.11,
\]
\[
    q(1,0)\approx 0.974,
    \qquad
    q(1,1)=-1+0.1(-1)=-1.1.
\]

So the sample-final answer is
\[
    \boxed{
    v(0)\approx -0.156,\quad
    v(1)\approx 0.974
    }
\]
and
\[
    \boxed{
    q(0,0)=-1,\quad
    q(0,1)=1.11,\quad
    q(1,0)\approx 0.974,\quad
    q(1,1)=-1.1
    }.
\]

\subsection*{Part (b): Off-policy TD control (Q-learning)}

Use
\[
    q(s_t,a_t)\leftarrow q(s_t,a_t)
    + \alpha\left[r_{t+1}+\gamma\max_{a'}q(s_{t+1},a')-q(s_t,a_t)\right].
\]

Start from
\[
    q(s,a)=0\quad \text{for all }(s,a).
\]

\subsubsection*{Update along \(\tau_0\)}

\[
    q(0,1)=0+0.6[1+0.1(0)-0]=0.6.
\]

\[
    q(1,0)=0+0.6[1+0.1(0)-0]=0.6.
\]

\[
    q(1,0)=0.6+0.6[1+0.1(0.6)-0.6]
    =0.6+0.6(0.46)=0.876.
\]

\subsubsection*{Update along \(\tau_1\)}

\[
    q(1,0)=0.876+0.6[1+0.1(0.876)-0.876]
    \approx 1.003.
\]

\[
    q(1,1)=0+0.6[-1+0.1(0.6)-0]=-0.564.
\]

\[
    q(0,0)=0+0.6[-1+0.1(0.6)-0]=-0.564.
\]

Therefore
\[
    \boxed{
    q(0,0)=-0.564,\quad
    q(0,1)=0.6,\quad
    q(1,0)\approx 1.003,\quad
    q(1,1)=-0.564
    }.
\]

\section*{Question 6: Bayesian Network}

Variables:
\begin{itemize}
    \item \(S\): storm,
    \item \(B\): burglar,
    \item \(C\): cat,
    \item \(A\): alarm.
\end{itemize}

\subsection*{Part (a): Directed model}

The natural causal graph is
\[
    S \to B,\qquad S \to C,\qquad B \to A,\qquad C \to A.
\]

Reason:
\begin{itemize}
    \item storms affect whether burglars show up,
    \item storms affect whether cats come inside,
    \item burglars and cats both affect whether the alarm is triggered.
\end{itemize}

\subsection*{Part (b): Estimate the probabilities}

The data table has 13 samples.

\subsubsection*{1. Prior for storm}

Storm is true in IDs \(10,11,12,13\), so
\[
    P(S=T)=\frac{4}{13}\approx 0.3077.
\]

\subsubsection*{2. Burglar given storm}

When \(S=T\), only ID \(13\) has \(B=T\), so
\[
    P(B=T\mid S=T)=\frac{1}{4}=0.25.
\]

When \(S=F\), IDs \(7,8,9\) have \(B=T\), so
\[
    P(B=T\mid S=F)=\frac{3}{9}=\frac13\approx 0.3333.
\]

\subsubsection*{3. Cat given storm}

When \(S=T\), IDs \(10,11,12\) have \(C=T\), so
\[
    P(C=T\mid S=T)=\frac{3}{4}=0.75.
\]

When \(S=F\), IDs \(6,9\) have \(C=T\), so
\[
    P(C=T\mid S=F)=\frac{2}{9}\approx 0.2222.
\]

\subsubsection*{4. Alarm given burglar and cat}

\[
    P(A=T\mid B=T,C=T)=\frac{1}{1}=1.0
\]
because ID \(9\) is the only case.

\[
    P(A=T\mid B=T,C=F)=\frac{2}{3}\approx 0.6667
\]
from IDs \(7,8,13\), where alarm is true in \(8,13\).

\[
    P(A=T\mid B=F,C=T)=\frac{1}{4}=0.25
\]
from IDs \(6,10,11,12\), where alarm is true only in \(10\).

\[
    P(A=T\mid B=F,C=F)=\frac{1}{5}=0.2
\]
from IDs \(1,2,3,4,5\), where alarm is true only in \(5\).

So the CPTs are:
\[
    \boxed{P(S=T)=0.3077}
\]
\[
    \boxed{
    P(B=T\mid S=T)=0.25,\quad
    P(B=T\mid S=F)=0.3333
    }
\]
\[
    \boxed{
    P(C=T\mid S=T)=0.75,\quad
    P(C=T\mid S=F)=0.2222
    }
\]
\[
    \boxed{
    \begin{aligned}
    P(A=T\mid B=T,C=T)&=1.0,\\
    P(A=T\mid B=T,C=F)&=0.6667,\\
    P(A=T\mid B=F,C=T)&=0.25,\\
    P(A=T\mid B=F,C=F)&=0.2.
    \end{aligned}
    }
\]

\subsection*{Part (c): Predict ALARM when \(B=T\), \(C=T\), \(S=F\)}

Once \(B\) and \(C\) are known, \(A\) depends only on \(B\) and \(C\), not directly
on \(S\).

From the CPT:
\[
    P(A=T\mid B=T,C=T)=1.
\]

Therefore the Bayesian network predicts
\[
    \boxed{A=T.}
\]

\subsection*{Part (d): Generate one sample}

Use the random numbers
\[
    0.94,\ 0.27,\ 0.30,\ 0.18.
\]

Following the same sampling convention as the sample solution, assign the
``false'' interval first and the ``true'' interval second.

\subsubsection*{1. Sample \(S\)}

\[
    P(S=T)=0.3077
    \quad\Rightarrow\quad
    P(S=F)=0.6923.
\]

Since
\[
    0.94 > 0.6923,
\]
we get
\[
    S=T.
\]

\subsubsection*{2. Sample \(B\mid S=T\)}

\[
    P(B=T\mid S=T)=0.25
    \quad\Rightarrow\quad
    P(B=F\mid S=T)=0.75.
\]

Since
\[
    0.27 < 0.75,
\]
we get
\[
    B=F.
\]

\subsubsection*{3. Sample \(C\mid S=T\)}

\[
    P(C=T\mid S=T)=0.75
    \quad\Rightarrow\quad
    P(C=F\mid S=T)=0.25.
\]

Since
\[
    0.30 > 0.25,
\]
we get
\[
    C=T.
\]

\subsubsection*{4. Sample \(A\mid B=F,C=T\)}

\[
    P(A=T\mid B=F,C=T)=0.25
    \quad\Rightarrow\quad
    P(A=F\mid B=F,C=T)=0.75.
\]

Since
\[
    0.18 < 0.75,
\]
we get
\[
    A=F.
\]

Therefore the generated sample is
\[
    \boxed{
    S=T,\quad B=F,\quad C=T,\quad A=F.
    }
\]

\section*{Final Answers Summary}

\begin{itemize}
    \item Q1:
    \[
        \hat y = 0.5648 + 0.2963x,\quad
        \text{MSE}_{\text{train}}\approx 2.0,\quad
        \text{MSE}_{\text{test}}\approx 0.206.
    \]
    \item Q2: use the provided softmax derivative chain to obtain
    \[
        \nabla_{\beta_k}J(\beta_k)
        =
        -\sum_{i=1}^{M}\sum_{k=1}^{K}
        I_k(y^{(i)})
        \left(1-\text{softmax}(z^{(i)})_k\right)x^{(i)}.
    \]
    \item Q3:
    \[
        \delta_1=0,\ \delta_2=-0.095,\ \delta_3=-0.2375,\ \delta_4=0,\ 
        \delta_5=-0.095,\ \delta_6=-0.2375,\ \delta_7=-0.475
    \]
    and
    \[
        \frac{\partial L}{\partial w_{75}}=-1.30625,\qquad
        \frac{\partial L}{\partial w_1}=-0.3325.
    \]
    \item Q4:
    \[
        G_0=65,\qquad G_1=70.
    \]
    \item Q5:
    \[
        v(0)\approx -0.156,\quad v(1)\approx 0.974
    \]
    and
    \[
        q(0,0)=-1,\ q(0,1)=1.11,\ q(1,0)\approx 0.974,\ q(1,1)=-1.1
    \]
    for the sample-solution Monte Carlo part, and
    \[
        q(0,0)=-0.564,\ q(0,1)=0.6,\ q(1,0)\approx 1.003,\ q(1,1)=-0.564
    \]
    for Q-learning.
    \item Q6:
    \[
        P(S=T)=0.3077,\quad
        P(B=T\mid S=T)=0.25,\quad
        P(B=T\mid S=F)=0.3333,
    \]
    \[
        P(C=T\mid S=T)=0.75,\quad
        P(C=T\mid S=F)=0.2222,
    \]
    \[
        P(A=T\mid T,T)=1,\quad
        P(A=T\mid T,F)=0.6667,\quad
        P(A=T\mid F,T)=0.25,\quad
        P(A=T\mid F,F)=0.2,
    \]
    with prediction \(A=T\) for \(B=T,C=T,S=F\), and generated sample
    \[
        (S,B,C,A)=(T,F,T,F).
    \]
\end{itemize}

\end{document}
